\section{Keeping the Season}

The Sundays and days give the season its texts. This part turns to the work they ask of the one
who keeps them: the watching the Fathers teach, the renewal of prayer, fasting and almsgiving,
a rule of life, and confession before Christmas.

\subsection{Watching with the Fathers}\label{sec:fathers}

The command under all the season's waiting is the Lord's own:

\begin{quote}
Watch ye therefore, because you know not what hour your Lord will come.\par
{\smallskip\noindent\small ---Matthew 24:42, Douay-Rheims (Challoner).\par}
\end{quote}

\noindent The passages gathered here stand apart from the Sunday Gospels. They
are what four Fathers and one Doctor taught about the waiting itself: how to watch, how to pray
while waiting, what to do with the word of God, how Mary waited, and what the waiting asks of the
hands.

\subsubsection{Watching without reckoning}

St Cyril of Jerusalem, in his fifteenth catechetical lecture, on the second coming of Christ,
warned his hearers off one kind of interest in it:

\begin{quote}
The things then which are seen shall pass away, and there shall come the things which are looked
for, things fairer than the present; but as to the time let no one be curious. For it is not for
you, He says, to know times or seasons, which the Father hath put in His own power. And venture
not thou to declare when these things shall be, nor on the other hand supinely slumber. For he
saith, Watch, for in such an hour as ye expect not the Son of Man cometh.\par
{\smallskip\noindent\small ---Cyril of Jerusalem, \textit{Catechesis} XV.4, tr. E.~H. Gifford,
NPNF, 2nd ser., vol.~7 (New York, 1894).\par}
\end{quote}

The Gospel of the First Sunday describes the signs, and Cyril keeps the reader from making a
calendar of them.

The things looked for are ``fairer than the present.'' Cyril does not set the coming of Christ
before his hearers as a threat only. It is the passing away of what is seen and the arrival of
what is hoped for. For the last things themselves, death, judgment, heaven and hell, which the
season's Gospels bring before the reader, see the essays on the Four Last Things in this library's
Faith collection.

\subsubsection{Desire that prays}

Watching has to go on through a working day, and St Augustine, writing in a letter about prayer,
found the means of it in desire:

\begin{quote}
When we cherish uninterrupted desire along with the exercise of faith and hope and charity, we
``pray always.'' But at certain stated hours and seasons we also use words in prayer to God, that
by these signs of things we may admonish ourselves, and may acquaint ourselves with the measure of
progress which we have made in this desire, and may more warmly excite ourselves to obtain an
increase of its strength. For the effect following upon prayer will be excellent in proportion to
the fervour of the desire which precedes its utterance.\par
{\smallskip\noindent\small ---Augustine, \textit{Epistula} CXXX.18, tr. J.~G. Cunningham, NPNF,
1st ser., vol.~1 (Buffalo, 1886).\par}
\end{quote}

Augustine goes on in the same letter to read ``Pray without ceasing'' as ``Desire without
intermission, from Him who alone can give it, a happy life, which no life can be but that which is
eternal.'' In his exposition of the
thirty-seventh Psalm (the thirty-eighth in the English numbering) he says the same thing more
plainly still: the desire of the heart is itself a prayer, and if the desire does not cease, the
prayer does not cease. In Advent the words said at stated hours, the Collect in the morning and
the O antiphon at night, are such signs, by which the soul admonishes itself and measures how far
its desire has grown.

\subsubsection{Keeping the word}

St Bernard, having named the hidden coming of Christ in the soul (Section~\ref{sec:three-comings}),
tells his monks how it is received. The answer is the word of God, kept as food is kept:

\begin{quote}
Keep the word of God in the same way in which the food of thy body can be best preserved; for it
is a living bread and the food of thy soul. Material bread, whilst it is kept in the cupboard may
be stolen by thieves, or nibbled by mice, or may become stale and unfit for use. But once it has
been assimilated, we have nothing further to fear from such possibilities. It is thus we should
keep the word of God, for ``blessed are they who keep it'' so. Receive this spiritual bread,
therefore, into the stomach of your minds, and having been there properly digested, let it be
transmitted thence to the affections and the will.\par
{\smallskip\noindent\small ---Bernard, \textit{In Adventu Domini} V.2, tr. a priest of Mount
Melleray (Dublin, [1921]).\par}
\end{quote}

\noindent Bernard goes on:

\begin{quote}
My brethren, if you thus keep the word of God, there can be no doubt that it, in turn, shall keep
you. For the Son will come to you with the Father, the great Prophet will come to you to rebuild
the Jerusalem of your souls and to make all things new.\par
{\smallskip\noindent\small ---Bernard, \textit{In Adventu Domini} V.3, tr. Mount Melleray ([1921]).\par}
\end{quote}

\noindent The Isaiah of the daily lessons, the Sunday Epistle and Gospel, or a sentence of a
Father can be kept in this way; Section~\ref{sec:prayer} describes one old method of doing so.

\subsubsection{With Mary}

No one waited for Christ as Mary did, and St Ambrose makes her a pattern for every believer. At
the Annunciation he notices her first word about herself. As the \textit{Catena Aurea} gives
his words:

\begin{quote}
Behold now the humility, the devotion of the virgin. For it follows, But Mary said, Behold the
handmaid of the Lord. She calls herself His handmaid, who is chosen to be His mother, so far was
she from being exalted by the sudden promise.\par
{\smallskip\noindent\small ---Ambrose, \textit{Expositio evangelii secundum Lucam} II, as the
\textit{Catena Aurea} gives it on Luke 1:38, Oxford tr. (1843).\par}
\end{quote}

At the Visitation, when Mary's soul magnifies the Lord, Ambrose turns her song outward to every
believer. Again as the \textit{Catena} gives it, in a condensed form of his commentary:

\begin{quote}
But let every one have the spirit of Mary, so that he may rejoice in the Lord. If according to the
flesh there is one mother of Christ, yet, according to faith, Christ is the fruit of all. For every
soul receives the word of God if only he be unspotted and free from sin, and preserves it with
unsullied purity.\par
{\smallskip\noindent\small ---Ambrose, \textit{Expositio evangelii secundum Lucam} II, as the
\textit{Catena Aurea} gives it on Luke 1:46--47, Oxford tr. (1843).\par}
\end{quote}

Set beside Bernard's sermon, Ambrose's sentence suggests a reading that is this book's own, not
either Father's: the soul that receives the word of God and keeps it bears Christ by faith, as
Mary bore him in the flesh.

\subsubsection{The lamp and the works}

Watching, praying and keeping the word are not idle. Cyril, at the end of the same lecture on the
second coming, asks how the judgment is to be met, and answers from the Gospel of the judgment:

\begin{quote}
I was an hungred, and ye gave Me meat; I was thirsty, and ye gave Me drink; I was a stranger, and ye
took Me in; naked, and ye clothed Me; I was sick, and ye visited Me; I was in prison, and ye came
unto Me. These things if thou do, thou shalt reign together with Him; but if thou do them not, thou
shalt be condemned. At once then begin to do these works, and abide in the faith; lest, like the
foolish virgins, tarrying to buy oil, thou be shut out. Be not confident because thou merely
possessest the lamp, but constantly keep it burning.\par
{\smallskip\noindent\small ---Cyril of Jerusalem, \textit{Catechesis} XV.26, tr. Gifford, NPNF,
2nd ser., vol.~7 (1894).\par}
\end{quote}

\noindent Cyril says just before that ``there is here no need of allegory, but to fulfil what is
said'': the works he names are to be done as he names them.

St Leo the Great, preaching the December fast, asks for a fast of the mind to go with the fast of
the body: ``When the body therefore fasts from food, let the mind fast from vices, and pass
judgment upon all earthly cares and desires according to the law of its King'' (\textit{Sermo}
XIX.2, tr. Feltoe). Section~\ref{sec:fasting} turns to the fast of the body and the alms that go
with it.
